\documentclass[11pt]{article}

\usepackage{amsmath,amssymb,amsfonts}
\usepackage{hyperref}
\usepackage{graphicx}
\usepackage{booktabs}
\usepackage{geometry}
\geometry{margin=1in}

\title{A Stochastic Structural Framework for Measuring\\
Implied Bitcoin Accretion in MicroStrategy Equity:\\
Corrected Model and Implementation Plan}
\author{}
\date{\today}

\begin{document}
\maketitle

\section{Objectives and Scope}

The goal of this project is to:
\begin{itemize}
  \item Build a mathematically consistent structural model linking Bitcoin (BTC) price dynamics, MicroStrategy (MSTR) BTC holdings, its balance-sheet leverage, and the equity premium over net asset value (NAV).
  \item Distinguish between:
  \begin{itemize}
    \item \emph{Structural balance-sheet leverage} (from assets and debt).
    \item \emph{Total equity elasticity} (including the response of the premium to BTC).
    \item \emph{Total BTC growth} versus \emph{per-share BTC accretion}.
  \end{itemize}
  \item Calibrate the model to publicly available data.
  \item Compute and visualize:
    \begin{itemize}
      \item Implied BTC Growth Rate (IBGR) -- total and per-share.
      \item Implied Leverage Elasticity (ILE).
      \item Total Equity Elasticity (TEE).
      \item Premium Mean-Reversion Index (PMRI).
      \item Implied Funding Requirement Distribution (IFRD).
      \item Survival / recapitalization probabilities.
    \end{itemize}
  \item Implement the full pipeline in Python with clear data ingestion, calibration, simulation, and reporting.
\end{itemize}

\section{Model Specification}

\subsection{State Variables, Measures, and Balance Sheet}

We work on a filtered probability space
$(\Omega,\mathcal F,(\mathcal F_t)_{t \ge 0})$
with two measures:
\begin{itemize}
  \item $\mathbb P$: historical (real-world) measure, used for premium dynamics, management behavior (holdings), and survival/IFRD analysis.
  \item $\mathbb Q$: risk-neutral measure, used (if needed) to calibrate BTC volatility/jump parameters to options.
\end{itemize}

State variables:
\begin{itemize}
  \item $S_t$: BTC (BTC-USD) price.
  \item $H_t$: MicroStrategy BTC holdings (in BTC).
  \item $D_t$: outstanding debt (face-value), modeled deterministically at baseline.
  \item $A_t$: BTC asset value,
  $$
    A_t = H_t S_t.
  $$
  \item $P_t$: MSTR share price.
  \item $N_t^{\text{sh}}$: number of MSTR shares outstanding.
  \item $E_t$: equity market value,
  $$
    E_t = P_t N_t^{\text{sh}}.
  $$
  \item $B_t$: BTC per share,
  $$
    B_t = \frac{H_t}{N_t^{\text{sh}}}.
  $$
\end{itemize}

Define net asset value (NAV) with limited liability:
$$
  \text{NAV}_t = (A_t - D_t)^+ = \max(H_t S_t - D_t, 0).
$$

When computing the historical log-premium from data we will:
\begin{itemize}
  \item Introduce a small floor $\varepsilon > 0$ to avoid log singularities:
  $$
    \text{NAV}_t^{\text{clip}} = \max(A_t - D_t, \varepsilon),
  $$
  \item and define
  $$
    \pi_t = \log\left(\frac{E_t}{\text{NAV}_t^{\text{clip}}}\right).
  $$
\end{itemize}
Periods where $A_t \approx D_t$ will be tagged as a ``distress regime'' and treated carefully in calibration.

\subsection{BTC Price Dynamics}

Under the risk-neutral measure $\mathbb Q$, for option calibration we consider a jump-diffusion:
$$
  \frac{dS_t}{S_{t^-}} =
  \left(r_t - q_S - \lambda_S k_S \right) dt
  + \sigma_S(S_t,t)\,dW_t^{S,\mathbb Q}
  + \left(e^{Y_S} - 1 \right) dN_t^S,
$$
where:
\begin{itemize}
  \item $r_t$ is the risk-free rate (deterministic or piecewise constant).
  \item $q_S$ is any BTC ``carry'' term (set $q_S = 0$ initially).
  \item $W_t^{S,\mathbb Q}$ is a Brownian motion under $\mathbb Q$.
  \item $N_t^S$ is a Poisson process with intensity $\lambda_S$.
  \item $Y_S$ is the (log-)jump size, with distribution $\nu_S$.
  \item $k_S = \mathbb E[e^{Y_S}-1]$ is the jump compensator.
  \item $\sigma_S(S_t,t)$ is a (possibly) state-dependent volatility surface calibrated to options.
\end{itemize}

For historical $\mathbb P$-measure simulations focusing on survival and IFRD, a simpler GBM with constant volatility fitted to BTC returns is acceptable as a first step:
$$
  \frac{dS_t}{S_t} = \mu_S dt + \sigma_S dW_t^{S,\mathbb P}.
$$

\subsection{Premium Dynamics under $\mathbb P$}

We model the log-premium $\pi_t$ as an Ornstein--Uhlenbeck (OU) process under $\mathbb P$ (management and sentiment world), with optional jumps:
$$
  d\pi_t = \kappa(\theta - \pi_t) dt
           + \sigma_\pi dW_t^{\pi,\mathbb P}
           + \eta_\pi dJ_t^\pi,
$$
where:
\begin{itemize}
  \item $\kappa > 0$ is the mean-reversion speed.
  \item $\theta$ is the long-run mean premium.
  \item $\sigma_\pi$ is the volatility of premium.
  \item $W_t^{\pi,\mathbb P}$ is a Brownian motion under $\mathbb P$.
  \item $J_t^\pi$ is a pure-jump process with intensity $\lambda_\pi$ (can be set to zero initially).
  \item The Brownian motions are correlated:
  $$
    \text{corr}(dW_t^{S,\mathbb P}, dW_t^{\pi,\mathbb P}) = \rho.
  $$
\end{itemize}

Discrete-time approximation for daily step $\Delta t$ (ignoring jumps initially):
$$
  \pi_{t+\Delta t}
  = \pi_t
    + \kappa(\theta - \pi_t) \Delta t
    + \sigma_\pi \sqrt{\Delta t}\,Z_\pi,
$$
with $(Z_S,Z_\pi)$ bivariate normal and $\text{corr}(Z_S,Z_\pi)=\rho$.

\subsection{BTC Holding Dynamics (Flywheel) under $\mathbb P$}

We encode the ``leverage flywheel'' via an endogenous BTC holding process:
$$
  dH_t = \alpha\,\pi_t^+ dt + dH_t^{\text{jump}},
  \quad \pi_t^+ = \max(\pi_t,0),
$$
where the jump component captures discrete financing events:
$$
  dH_t^{\text{jump}} = \sum_{k} \Delta H_{\tau_k}\,\delta_{t=\tau_k},
  \qquad \Delta H_{\tau_k} \ge 0.
$$

Here:
\begin{itemize}
  \item $\{\tau_k\}$ are event times of a Poisson process $M_t$ with intensity $\lambda_M$.
  \item $\Delta H_{\tau_k}$ are BTC amounts purchased using financing proceeds.
\end{itemize}

In discrete time:
\begin{itemize}
  \item Draw $N_t^M \sim \text{Poisson}(\lambda_M \Delta t)$.
  \item Draw jump sizes $\Delta H_{t,1},\dots,\Delta H_{t,N_t^M}$ from a distribution $\mathcal D_H$.
  \item Update
  $$
    H_{t+\Delta t}
    = H_t + \alpha \pi_t^+ \Delta t
          + \sum_{i=1}^{N_t^M} \Delta H_{t,i},
    \qquad H_{t+\Delta t} \ge 0.
  $$
\end{itemize}

Calibration choices:
\begin{itemize}
  \item $\alpha$ is estimated on aggregated (weekly/monthly) holding changes to avoid daily noise.
  \item If data suggests $\alpha \approx 0$, we can set $\alpha=0$ and model all accretion as jumps, potentially with $\lambda_M$ depending on the premium (e.g. $\lambda_M(\pi_t)$).
  \item As an extension, the jump intensity or sizes can be made to depend on leverage (e.g. $\beta_{LS}(t)$) to capture over-leverage constraints.
\end{itemize}

\subsection{Debt Process and Convertibles}

At baseline, we treat total debt as deterministic piecewise-constant:
$$
  D_t = \sum_{j=1}^J D_j \mathbf 1_{\{t < T_j\}},
$$
where each instrument $j$ has face value $D_j$ and maturity $T_j$.
This is conservative for solvency analysis, especially for convertible notes (we ignore their equity-like behavior when deep in the money).

In implementation, each instrument will be represented by a \texttt{DebtInstrument} object with fields:
\begin{itemize}
  \item Face value, coupon, issue date, maturity date.
  \item Convertible flag, conversion price, and other relevant terms.
\end{itemize}

This allows an extension where an ``effective debt'' $D_t^{\text{eff}}$ discounts the equity-like component of deep in-the-money convertibles.

\subsection{Share-Count Dynamics and BTC per Share}

To distinguish total BTC growth from per-share accretion, we model the share count:
$$
  dN_t^{\text{sh}} = dN_t^{\text{(equity)}} + dN_t^{\text{(conv)}},
$$
where:
\begin{itemize}
  \item $dN_t^{\text{(equity)}}$ captures new stock issuance (e.g. ATMs, follow-ons).
  \item $dN_t^{\text{(conv)}}$ captures conversion of convertible notes.
\end{itemize}

In practice, both components are lumpy and are best modeled as jump processes with event times aligned with financing and conversion events.

The BTC per share process is:
$$
  B_t = \frac{H_t}{N_t^{\text{sh}}}.
$$
Per-share accretion metrics will be defined in terms of $B_t$.

\subsection{Equity Valuation, Structural Leverage, and Total Elasticity}

The structural relation between equity, assets, debt, and premium is:
$$
  E_t = e^{\pi_t} \cdot (A_t - D_t)^+ = e^{\pi_t} (H_t S_t - D_t)^+.
$$

\paragraph{Implied Leverage Elasticity (ILE).}
Define the \emph{structural} leverage elasticity as the sensitivity of equity to BTC price, ignoring instantaneous changes in premium:
$$
  \beta_{LS}(t)
  = \frac{\partial \log E_t}{\partial \log S_t}\bigg|_{\pi_t \text{ fixed}}
  \approx \frac{H_t S_t}{H_t S_t - D_t},
$$
whenever $H_t S_t > D_t$.
This depends only on the current balance sheet and is independent of $\pi_t$.

\paragraph{Total Equity Elasticity (TEE).}
In reality, $\pi_t$ typically depends on BTC shocks. From
$$
  \log E_t = \pi_t + \log(H_t S_t - D_t),
$$
one obtains
$$
  \frac{d \log E_t}{d \log S_t}
  = \frac{d \pi_t}{d \log S_t}
    + \frac{H_t S_t}{H_t S_t - D_t}.
$$
Define the \emph{Total Equity Elasticity}
$$
  \beta_{\text{TOT}}(t)
  = \frac{d \log E_t}{d \log S_t}
  \approx \gamma_{\pi S} + \beta_{LS}(t),
$$
where $\gamma_{\pi S}$ is the empirical sensitivity
$$
  \gamma_{\pi S} \approx
  \frac{\text{Cov}(\Delta \pi_t, \Delta \log S_t)}{\text{Var}(\Delta \log S_t)},
$$
estimated via regression.
We will report both $\beta_{LS}(t)$ (pure balance-sheet leverage) and $\beta_{\text{TOT}}(t)$ (what traders \emph{feel} in equity P\&L).

\section{Implied Quantitative Indicators}

\subsection{Implied BTC Growth Rate (IBGR)}

\paragraph{Total IBGR.}
Under the historical measure $\mathbb P$, define the expected total BTC holding growth over a short horizon $\Delta t$:
$$
  \mu_H
  = \frac{1}{H_0}
    \mathbb E^\mathbb P
    \left[
      \frac{H_{\Delta t} - H_0}{\Delta t}
      \,\bigg|\,
      S_0, \pi_0, H_0
    \right].
$$
Using the holding dynamics for small $\Delta t$,
$$
  \mu_H \approx
  \frac{1}{H_0}
  \left(
    \alpha \pi_0^+ + \lambda_M \mathbb E[\Delta H]
  \right),
$$
where $\mathbb E[\Delta H]$ is the mean BTC added per financing event.

\paragraph{Per-Share IBGR.}
Define BTC per share $B_t = H_t / N_t^{\text{sh}}$.
Per-share accretion is captured by
$$
  \mu_{H/N}
  =
  \frac{1}{B_0}
  \mathbb E^\mathbb P
  \left[
    \frac{B_{\Delta t} - B_0}{\Delta t}
    \,\bigg|\,
    S_0, \pi_0, H_0, N_0^{\text{sh}}
  \right].
$$
Exact closed form is messy; in practice we will:
\begin{itemize}
  \item Simulate joint paths of $(H_t,N_t^{\text{sh}})$ under the calibrated dynamics.
  \item Estimate $\mu_{H/N}$ via Monte Carlo.
\end{itemize}
Positive $\mu_{H/N}$ indicates \emph{accretive} BTC growth per share.

\subsection{Implied Leverage Elasticity (ILE)}

As above, for a given state $(S_t,H_t,D_t)$,
$$
  \text{ILE}_t = \beta_{LS}(t)
  = \frac{H_t S_t}{H_t S_t - D_t},
$$
whenever $H_t S_t > D_t$.
This is purely a function of balance-sheet leverage.

\subsection{Total Equity Elasticity (TEE)}

Given the empirical $\gamma_{\pi S}$ and current $\beta_{LS}(t)$:
$$
  \beta_{\text{TOT}}(t)
  \approx \gamma_{\pi S} + \beta_{LS}(t).
$$
This indicator captures the full sensitivity of equity to BTC, including premium dynamics.

\subsection{Premium Mean-Reversion Index (PMRI)}

For the OU process (ignoring jumps), the stationary variance is
$$
  \text{Var}^{\text{ss}}(\pi_t) = \frac{\sigma_\pi^2}{2\kappa}.
$$
Define the z-score:
$$
  \text{PMRI} =
  \frac{\pi_0 - \theta}{\sigma_\pi / \sqrt{2\kappa}},
$$
with large positive values indicating exuberant premium, and negative values indicating discount/stress.

\subsection{Implied Funding Requirement Distribution (IFRD)}

The cumulative dollar cost of BTC acquisitions over horizon $\Delta t$ is
$$
  G_{\Delta t} = \int_0^{\Delta t} S_u\,dH_u.
$$
Using the decomposition:
$$
  G_{\Delta t}
  = \int_0^{\Delta t} S_u \alpha \pi_u^+\,du
    + \sum_{0 < u \le \Delta t} S_u \Delta H_u.
$$

In discrete form for simulated paths:
$$
  G_{\Delta t}^{(\text{path})}
  \approx \sum_{k=0}^{n-1} S_{t_k} (H_{t_{k+1}} - H_{t_k}).
$$
Monte Carlo over many paths provides the empirical distribution of $G_{\Delta t}$:
mean, variance, and VaR/ES (e.g. $\text{VaR}_{95\%}(G_{\Delta t})$).

\subsection{Survival and Recapitalization Probabilities}

Define a distress threshold with buffer $\varepsilon \ge 0$:
$$
  A_t \le (1+\varepsilon) D_t
  \quad \Longleftrightarrow \quad
  H_t S_t \le (1+\varepsilon) D_t.
$$
Let
$$
  \tau = \inf \{ u \in [0,T] : A_u \le (1+\varepsilon) D_u \}
$$
be the first-passage time to distress.

Simulating under $\mathbb P$:
\begin{itemize}
  \item The survival probability is
  $$
    \mathbb P(\tau > T) \approx 
    \frac{\#\{\text{paths with no hit before } T\}}{\#\{\text{total paths}\}}.
  $$
  \item Recapitalization probabilities can be defined with thresholds on equity, premium, or funding requirements.
\end{itemize}

\section{Data Requirements and Download Instructions}

All series should be aligned to a common daily date index.

\subsection{BTC Price and Derivatives}

\paragraph{BTC Spot/OHLCV.}
\begin{itemize}
  \item BTC-USD daily close (or OHLCV) series $S_t$.
  \item Sources:
  \begin{itemize}
    \item Yahoo Finance (\texttt{"BTC-USD"}), via web or \texttt{yfinance}.
    \item CryptoDataDownload / CoinGecko / CoinMarketCap for CSV downloads.
  \end{itemize}
\end{itemize}

\paragraph{BTC Options (Optional).}
For a richer volatility/jump calibration under $\mathbb Q$:
\begin{itemize}
  \item Deribit BTC options chain.
  \item CME BTC futures/options from a data vendor.
\end{itemize}

\subsection{MSTR Equity Data}

\begin{itemize}
  \item Daily MSTR share price $P_t$.
  \item Shares outstanding $N_t^{\text{sh}}$ (quarterly + event-based).
  \item Sources:
  \begin{itemize}
    \item Yahoo Finance (\texttt{"MSTR"}) for price.
    \item SEC filings (10-K, 10-Q), MSTR investor relations, data sites for share count history.
  \end{itemize}
\end{itemize}

\subsection{MicroStrategy BTC Holdings}

\begin{itemize}
  \item BTC holdings time series $H_t$ (in BTC).
  \item Purchase events table: date, BTC added, USD notional, average price.
  \item Sources:
  \begin{itemize}
    \item MicroStrategy / company site disclosures.
    \item Corporate BTC tracking sites with CSV downloads.
  \end{itemize}
\end{itemize}

\subsection{Debt Schedule and Financing Events}

\begin{itemize}
  \item Debt instruments: issue date, maturity, face value, coupon, convertible terms.
  \item Financing events: equity issuance, new debt, conversions.
  \item Sources:
  \begin{itemize}
    \item SEC EDGAR (10-K, 10-Q, 8-K).
    \item Company press releases / investor presentations.
  \end{itemize}
\end{itemize}

\subsection{Risk-Free Curve}

\begin{itemize}
  \item US Treasury yields (e.g. 1Y, 2Y) from FRED.
  \item Simple version: use a constant annual risk-free rate (e.g. 3\%).
\end{itemize}

\section{Estimation and Calibration Plan}

\subsection{Data Ingestion and Alignment}

\begin{enumerate}
  \item Download BTC, MSTR price series, holdings, debt schedule, shares outstanding as CSV or via APIs.
  \item In Python (\texttt{pandas}), build a daily date index from first BTC purchase to latest data.
  \item Merge:
  \begin{itemize}
    \item Forward-fill $H_t$, $D_t$, and $N_t^{\text{sh}}$ between known dates.
    \item For non-trading days for MSTR, forward-fill $P_t$ or work on trading-day-only calendar.
  \end{itemize}
  \item Compute derived series:
  \begin{itemize}
    \item $A_t = H_t S_t$.
    \item $\text{NAV}_t = (A_t - D_t)^+$.
    \item $\text{NAV}_t^{\text{clip}} = \max(A_t - D_t,\varepsilon)$.
    \item $E_t = P_t N_t^{\text{sh}}$.
    \item $\pi_t = \log(E_t / \text{NAV}_t^{\text{clip}})$.
    \item $B_t = H_t / N_t^{\text{sh}}$.
  \end{itemize}
  \item Tag dates with $A_t \approx D_t$ (or $\text{NAV}_t$ close to zero) as distress regime indicators.
\end{enumerate}

\subsection{Calibration of Premium OU under $\mathbb P$}

\begin{enumerate}
  \item Use daily (or business-day) steps $\Delta t = 1/252$.
  \item Exclude distress-regime datapoints or treat them separately.
  \item Use the exact OU transition:
  $$
    \pi_{t+\Delta t}
    = \theta + (\pi_t - \theta)e^{-\kappa \Delta t}
      + \sigma_\pi \sqrt{\frac{1 - e^{-2\kappa\Delta t}}{2\kappa}}\,Z_t,
  $$
  and perform maximum likelihood or nonlinear least squares to estimate $(\kappa,\theta,\sigma_\pi)$.
  \item Estimate $\rho$ by empirical correlation between BTC log returns $\Delta \log S_t$ and premium changes $\Delta \pi_t$.
  \item Estimate $\gamma_{\pi S}$ via regression of $\Delta \pi_t$ on $\Delta \log S_t$ to obtain total elasticity.
\end{enumerate}

\subsection{Calibration of BTC Holding Dynamics}

\begin{enumerate}
  \item Construct $\Delta H_t$ (changes in holdings) and identify financing event dates from disclosures.
  \item Build weekly or monthly aggregates:
  $$
    \Delta H^{(W)}_t, \quad \overline{\pi_t^+}^{(W)}.
  $$
  \item On non-event periods, regress:
  $$
    \Delta H^{(W)}_t \approx \alpha \overline{\pi_t^+}^{(W)} \Delta t + \varepsilon_t,
  $$
  to estimate $\alpha$ (if statistically significant).
  \item Treat event days as jumps:
  \begin{itemize}
    \item Jump sizes are the $\Delta H_t$ on those dates.
    \item Jump intensity $\lambda_M \approx N_{\text{events}} / T_{\text{years}}$.
    \item Estimate $\mathbb E[\Delta H]$ and $\mathcal D_H$ (empirical distribution).
  \end{itemize}
  \item If $\alpha$ is not significantly different from zero, set $\alpha=0$ and move to a pure jump-based accumulation model (possibly with $\lambda_M$ made premium-dependent).
\end{enumerate}

\subsection{Calibration of BTC Price Dynamics}

\begin{itemize}
  \item Baseline: estimate constant volatility $\sigma_S$ from historical BTC log returns (under $\mathbb P$).
  \item If using options, calibrate more sophisticated $\sigma_S(S_t,t)$ and jump parameters $(\lambda_S,\nu_S)$ under $\mathbb Q$.
\end{itemize}

\subsection{Share-Count and Convertible Calibration}

\begin{itemize}
  \item From equity and convertible issuance disclosures, reconstruct
        $N_t^{\text{sh}}$ and event-driven changes.
  \item Optionally model future issuance/convertible conversion behavior as a function of premium and leverage (e.g. intensities depending on $\pi_t$ and $\beta_{LS}(t)$).
\end{itemize}

\subsection{Indicator Computation}

\begin{enumerate}
  \item Use closed-form formulas where available:
  \begin{itemize}
    \item IBGR (total) from $(\alpha,\lambda_M,\mathbb E[\Delta H])$ and $H_0$.
    \item ILE from $(H_t,S_t,D_t)$.
    \item PMRI from OU parameters and current $\pi_0$.
    \item TEE from $\gamma_{\pi S}$ and current $\beta_{LS}(t)$.
  \end{itemize}
  \item For IFRD, per-share IBGR, and survival probabilities, run Monte Carlo simulations under $\mathbb P$.
\end{enumerate}

\section{Implementation Plan in Python}

\subsection{Project Structure}

A possible directory layout:
\begin{verbatim}
project_root/
  data_raw/
  data_processed/
  src/
    data_io.py
    preprocessing.py
    instruments.py
    calibration.py
    simulation.py
    indicators.py
    plots.py
  notebooks/
  results/
\end{verbatim}

\subsection{Core Modules}

\paragraph{1. \texttt{data\_io.py}}
\begin{itemize}
  \item \texttt{load\_btc\_prices():} load/download BTC-USD historical prices.
  \item \texttt{load\_mstr\_prices():} load/download MSTR prices.
  \item \texttt{load\_btc\_holdings():} read BTC holdings and event history.
  \item \texttt{load\_debt\_schedule():} read debt instruments.
  \item \texttt{load\_shares\_outstanding():} reconstruct $N_t^{\text{sh}}$.
\end{itemize}

\paragraph{2. \texttt{instruments.py}}
\begin{itemize}
  \item Define a \texttt{DebtInstrument} class with maturity, coupon, face value, and convertible fields.
  \item Provide methods to compute outstanding face at time $t$ and, later, effective debt.
\end{itemize}

\paragraph{3. \texttt{preprocessing.py}}
\begin{itemize}
  \item \texttt{build\_daily\_panel():}
  \begin{itemize}
    \item Merge BTC, MSTR, holdings, debt, and shares.
    \item Compute $A_t$, $\text{NAV}_t$, $\text{NAV}_t^{\text{clip}}$, $E_t$, $\pi_t$, $B_t$.
    \item Tag distress regime indicators.
  \end{itemize}
\end{itemize}

\paragraph{4. \texttt{calibration.py}}
\begin{itemize}
  \item \texttt{fit\_ou\_premium(pi\_series):} estimate $(\kappa,\theta,\sigma_\pi)$ and OU diagnostics.
  \item \texttt{estimate\_rho\_and\_gamma(btc\_returns, premium\_changes):} estimate $\rho$ and $\gamma_{\pi S}$.
  \item \texttt{fit\_holding\_dynamics(H\_series, pi\_series, event\_flags):} estimate $\alpha,\lambda_M,\mathbb E[\Delta H]$.
  \item \texttt{fit\_btc\_vol(btc\_returns):} estimate $\sigma_S$.
\end{itemize}

\paragraph{5. \texttt{simulation.py}}
\begin{itemize}
  \item Vectorized simulation of $(S_t,\pi_t,H_t,N_t^{\text{sh}},D_t,E_t)$ under $\mathbb P$.
  \item Implementation details:
  \begin{itemize}
    \item Construct the correlation matrix
    $$
      \Sigma =
      \begin{pmatrix}
        1 & \rho \\
        \rho & 1
      \end{pmatrix}
    $$
    and its Cholesky factor once.
    \item At each time step, draw independent standard normals and multiply by the Cholesky factor to obtain correlated $(Z_S,Z_\pi)$.
    \item Simulate jumps via Poisson draws and jump distributions.
  \end{itemize}
\end{itemize}

\paragraph{6. \texttt{indicators.py}}
\begin{itemize}
  \item \texttt{compute\_IBGR(params, state0):} compute total IBGR.
  \item \texttt{compute\_IBGR\_per\_share(simulated\_B\_paths):} estimate $\mu_{H/N}$ from Monte Carlo.
  \item \texttt{compute\_ILE(panel):} compute $\beta_{LS}(t)$ over time.
  \item \texttt{compute\_TEE(panel, gamma\_piS):} compute $\beta_{\text{TOT}}(t)$.
  \item \texttt{compute\_PMRI(params, pi0):} compute PMRI.
  \item \texttt{compute\_IFRD(sim\_paths):} compute IFRD distribution from pathwise $G_{\Delta t}$.
  \item \texttt{compute\_survival(sim\_paths, eps, T):} compute survival probabilities.
\end{itemize}

\paragraph{7. \texttt{plots.py}}
\begin{itemize}
  \item Time series plots for $(S_t,H_t,N_t^{\text{sh}},B_t,\pi_t,\text{NAV}_t,E_t)$.
  \item Scatter of $\Delta B_t$ vs $\pi_t$ (accretion vs premium).
  \item IFRD histograms, leverage elasticity distributions, and survival probability curves.
\end{itemize}

\subsection{Coding Checklist}

\begin{enumerate}
  \item Implement data loaders and build the daily panel.
  \item Fit OU premium parameters and estimate $\rho,\gamma_{\pi S}$.
  \item Estimate holding dynamics $(\alpha,\lambda_M,\mathbb E[\Delta H])$; decide if continuous term is necessary.
  \item Estimate BTC volatility and (optionally) jump parameters.
  \item Implement vectorized simulation of joint paths under $\mathbb P$.
  \item Compute IBGR (total and per-share), ILE, TEE, PMRI.
  \item Simulate IFRD and survival curves and generate plots.
\end{enumerate}

\section{Deliverables}

\begin{itemize}
  \item A cleaned, reproducible daily panel with $S_t,P_t,H_t,D_t,N_t^{\text{sh}},A_t,\text{NAV}_t,\text{NAV}_t^{\text{clip}},E_t,\pi_t,B_t$.
  \item A parameter report: $(\kappa,\theta,\sigma_\pi,\rho,\gamma_{\pi S},\alpha,\lambda_M,\mathbb E[\Delta H],\sigma_S)$.
  \item A Monte Carlo engine for $(S_t,\pi_t,H_t,N_t^{\text{sh}},D_t,E_t)$ under $\mathbb P$.
  \item Indicator dashboards (notebooks/scripts) presenting:
  \begin{itemize}
    \item Total and per-share IBGR.
    \item ILE time series and current value.
    \item TEE time series and current value.
    \item PMRI and its history.
    \item IFRD distribution and summary stats.
    \item Survival probabilities vs horizon $T$ and stress buffer $\varepsilon$.
    \item BTC-per-share vs premium plots to visually test accretion efficiency.
  \end{itemize}
  \item A written interpretation explaining whether current MSTR pricing implies realistic BTC accretion and leverage profiles, given historical behavior and calibrated parameters.
\end{itemize}

\end{document}